\paragraph{Controller cost}
EcoFed scores $|\mathcal C|=16$ configurations per client and sorts the best scores. Fig.~\ref{fig:scaling}a reports the measured time of the selection step in our (unoptimised, pure Python) implementation on synthetic fleets with 5\% of clients selected: 3~ms for $N=100$, 0.036~s for $N=10^3$, 0.26~s for $N=10^4$ and 2.7~s for $N=10^5$, 2.6 times the cost of an Oort-style top-$m$ ranking at $N=10^5$ (1.0~s) and 8 to 14 times at $N\le10^4$. This is negligible compared with a federated round, which takes seconds to minutes at that scale, but it is a Python measurement of the selection step only and says nothing about the scalability of learning itself, which we did not test beyond 20 clients.

\subsection{Cost model versus measured latency}\label{sec:latency}
The relative cost of depth is the lever behind depth control, so we checked it against wall-clock time. We timed one training step of the multi-exit network at depths 1--4 on a single CPU thread (batch 32, median of 40 repetitions). Measured time per sample grows more slowly than the modelled cost: depth 4 takes 3.4 times as long as depth 1 on UCI HAR's 9-channel input and 3.1 times on PAMAP2's 18-channel input, while the MAC count grows by a factor of 8.1 and 4.6 (Fig.~\ref{fig:scaling}b). A linear model of time in MACs fits well ($R^2=0.96$ and $0.98$) with a fixed cost per sample equivalent to about 0.75 and 0.55 million MACs. We therefore do not claim that the modelled savings from shallow updates are what a CPU would show; the sensitivity experiment of Section~\ref{sec:sens} injects that overhead into the true energy and shows that EcoFed's advantage on PAMAP2 shrinks but persists. Latency is not energy, and an accelerator would behave differently; this check only validates the \emph{relative} depth scaling on one CPU.

\begin{figure}[t]
\centering\includegraphics[width=\textwidth]{fig11_system}
\caption{(a) Time of the selection step against the number of clients (5\% selected; synthetic fleets; one CPU core). (b) Measured CPU time per training sample against the modelled cost, both relative to depth 1; the dashed line is proportionality.}\label{fig:scaling}
\end{figure}
